% !TeX root = ../main.tex
% This file is included by main.tex. Edit results.tex directly.

\section{Results}
\subsection{Overview of the evidence-guided wet-lab refinement workflow}


\begin{figure}[p]
  \centering
  \sffamily
  % \captionsetup{font=small, skip=3pt}

  \begin{adjustbox}{max width=\linewidth, max height=0.78\textheight, keepaspectratio}
  \begin{minipage}{\linewidth}
    \centering

    \begin{overpic}[width=0.96\linewidth]{Figure/Fig1/Flowchart_Figure_v3_a.pdf}
      \put(1,18){{\sffamily\bfseries a}}
    \end{overpic}

    \vspace{0.04in}

    \begin{overpic}[width=0.96\linewidth]{Figure/Fig1/Flowchart_Figure_v3_b.pdf}
      \put(1,25){{\sffamily\bfseries b}}
    \end{overpic}

    \vspace{0.04in}

    \begin{overpic}[width=0.96\linewidth]{Figure/Fig1/Flowchart_Figure_v3_c.pdf}
      \put(1,25){{\sffamily\bfseries c}}
    \end{overpic}

    \vspace{0.04in}

    \begin{overpic}[width=0.96\linewidth]{Figure/Fig1/Flowchart_Figure_v3_d.pdf}
      \put(1,40){{\sffamily\bfseries d}}
    \end{overpic}

  \end{minipage}
  \end{adjustbox}
  
    \caption{
    \textbf{Evidence-guided wet-lab refinement workflow.}
    \textbf{a}, Human-defined core agents and selected domain agents form a multi-role agent team.
    \textbf{b}, Raw literature is extracted into structured evidence.
    \textbf{c}, In silico tools generate sequence-level priors for planning.
    \textbf{d}, Evidence and priors enter team meetings to produce run-sheets, which are executed in wet-lab experiments and iteratively refined until a final recommendation is reached.
    }
  \label{fig:workflow_overview}
\end{figure}

We organized the optimization of $\beta$-LG glycation for reduced IgG-binding immunoreactivity as an evidence-guided, wet-lab-in-the-loop workflow linking a human researcher with a multi-role agent team (Fig.~\ref{fig:workflow_overview}). The human researcher defined the objective, experimental constraints and feasibility boundaries, while the agents structured literature evidence, generated sequence-level priors, formulated decision-focused run-sheets and interpreted returned results through comparator-aware discussion. These components formed a traceable evidence-to-experiment cycle in which each wet-lab round updated the questions and conditions prioritized in the next round, ultimately yielding a pre-specified recommendation for focused validation.

\subsection{Prior evidence and integrated in silico guidance for experimental decision-making}


\begin{figure}[p]
  \centering
  \sffamily
  % \captionsetup{font=small, skip=3pt}

  \begin{adjustbox}{max width=\linewidth, max height=0.90\textheight, keepaspectratio}
  \begin{minipage}{\linewidth}
    \centering

    \begin{overpic}[width=0.96\linewidth]{Figure/Fig2/Figure_2a.pdf}
      \put(1,17){{\sffamily\bfseries a}}
    \end{overpic}

    \vspace{0.04in}

    \begin{overpic}[width=0.96\linewidth]{Figure/Fig2/Figure_2b.pdf}
      \put(1,20){{\sffamily\bfseries b}}
    \end{overpic}

    \vspace{0.04in}

    \begin{overpic}[width=0.96\linewidth]{Figure/Fig2/Figure_2c.pdf}
      \put(1,40){{\sffamily\bfseries c}}
    \end{overpic}

    \vspace{0.04in}

    \begin{overpic}[width=0.96\linewidth]{Figure/Fig2/Figure_2d.pdf}
      \put(1,40){{\sffamily\bfseries d}}
    \end{overpic}

  \end{minipage}
  \end{adjustbox}
  
  \caption{
    \textbf{Structured literature evidence and sequence-level in silico priors for experimental decision-making.}
    \textbf{a}, Example extracted records from the structured evidence layer.
    \textbf{b}, Statistical overview of the literature-derived evidence layer.
    \textbf{c}, BepiPred-derived sequence-level epitope prior.
    \textbf{d}, NetGlycate-predicted glycation-prone lysine sites in $\beta$-LG.
  }
  \label{fig:evidence_prior}
\end{figure}

To establish an experimental evidence layer before wet-lab exploration, we assembled a project corpus of 23 full-text studies reporting glycation-related experimental conditions and immunoreactivity-related outcomes in $\beta$-LG or related food-protein systems. Using the Evidence Curator agent, these studies were parsed into standardized JSONL records that preserved experiment-specific information on reaction conditions, pretreatment settings, comparator relationships, immunoreactivity-related readouts, and reduction-related evidence. The extraction schema and controlled vocabulary are summarized in Supplementary Table~1, and an example excerpt from the structured evidence layer is shown in Fig.~\ref{fig:evidence_prior}a.

This process yielded a structured evidence base comprising 350 metric-level records from 207 unique experimental runs, including 280 treatment/comparison records and 70 comparator/control records. Because each JSONL row represents one experimental run measured by one allergenicity-related metric, a single experimental condition could contribute multiple records when multiple readouts were reported. The resulting evidence layer spanned the major decision axes of the optimization problem, including donor identity, reaction mode, temperature, reaction time, pretreatment, protein:sugar ratio, and state conditions such as pH, relative humidity, or water activity. Within this space, donor coverage was broad but uneven. Maltose and glucose were the most frequently recurring donor labels in the full structured dataset, while other repeatedly observed donors included galactose-related donors, mannose, lactose, ribose, and galacto-oligosaccharides. Among records with annotated reaction mode, dry-state reactions were more common than wet-state reactions (190 versus 79 records; 179 versus 73 among treatment/comparison records). Pretreatment coverage also showed an imbalance. At the full record level, records without pretreatment were more frequent than those with pretreatment (232 versus 118). When comparator/control records were excluded, no-pretreatment treatment/comparison records remained slightly more frequent than pretreatment records (164 versus 116), indicating that the imbalance was not solely driven by untreated controls. Temperature coverage was concentrated mainly in the 40--80~$^\circ$C range, whereas reaction time showed two dominant clusters, one at shorter durations of 2--8~h and another at longer durations of 24--72~h. The overall composition, field coverage, and major donor distribution of the literature-derived evidence layer are summarized in Fig.~\ref{fig:evidence_prior}b. Taken together, these records defined a broad but nonuniform experimental variable space, in which donor identities were numerous but repeated validation for any single donor was generally sparse.
The complete machine-readable evidence records are available in the project data repository.

At the same time, the structured evidence layer made the heterogeneity of the literature explicit. The most important sources of incompatibility were not simply differences in assay naming, but differences in condition backbones and comparator definitions. Across the extracted studies, comparator schemes varied across native or untreated baselines, heated or backbone controls, and treated-to-treated contrasts. As a result, many apparently favourable observations remained tied to local experimental backbones and could not be directly ranked across studies. This limitation was also reflected in the effect data themselves: among the 280 treatment/comparison records, only 68 originally contained directly reported numeric reduction percentages. To improve comparability, we standardized study outcomes into comparator-relative, direction-aware effect representations and linked each treatment/comparison record to its matched comparator whenever possible. After this harmonization, 154 treatment/comparison records could be represented with reduction-oriented effect information suitable for downstream comparison. Thus, the immediate value of the literature was not that it directly yielded a single executable optimal parameter set, but that it transformed a fragmented body of prior observations into a structured and substantially more comparable evidence layer for subsequent prioritization, filtering, and narrowing of the initial search space.

To complement the literature-derived evidence layer, the workflow next introduced sequence-level in silico guidance through a dedicated tool-selection and implementation stage. In team meetings involving the domain scientist roles, the Principal Investigator (PI), and the critic, the agent team considered which computational tools could provide experimentally relevant priors for the specific task of $\beta$-LG glycation optimization. The tool-screening rationale is summarized in Supplementary Table~2. On the basis of these discussions, two sequence-based tools were selected to provide complementary perspectives: one aimed at identifying residues more likely to be involved in B-cell epitope recognition, and the other aimed at identifying lysine sites more susceptible to non-enzymatic glycation. The human researcher then verified this tool choice and, under agent guidance, executed the selected analyses on the sequence of $\beta$-LG variant A. Importantly, the outputs of this stage were not treated as mechanistic proof or as hard constraints on subsequent decisions. Instead, they were incorporated together with the structured literature evidence as an additional layer of context for the next round of project specification and experimental prioritization.

The resulting sequence-level priors refined the early decision space in several ways. BepiPred-3.0~\cite{CHD+2022} identified higher-scoring putative epitope residues in several localized regions, with notable concentration around residues 61--69, 110--115, and 130--148, and with Lys69 and Lys77 also appearing among the top-ranked epitope-associated positions (Fig.~\ref{fig:evidence_prior}c). NetGlycate~\cite{JKB2006}, by contrast, assigned the highest glycation susceptibility scores to Lys47, Lys70, Lys14, Lys100, and Lys101, suggesting that these residues were among the most likely early sites of non-enzymatic modification (Fig.~\ref{fig:evidence_prior}d). Taken together, these outputs did not define a single mechanistically validated target region, but they did provide a more informative sequence context in which putative glycation-sensitive lysines could be interpreted alongside residues more likely to participate in immune recognition. In particular, the local proximity between the high-scoring glycation site at Lys70 and the epitope-associated region around residues 61--69 suggested that this part of the sequence warranted earlier attention during downstream condition prioritization. More broadly, the role of this in silico stage was not to replace wet-lab experimentation, but to increase the specificity of early-stage decision-making by supplying sequence-level priors that could be considered jointly with the extracted literature evidence.
Complete residue-level and lysine-level outputs are available in the project data repository and in the Source Data for Fig.~\ref{fig:evidence_prior}c,d.

Taken together, the structured literature evidence and the sequence-level in silico outputs shown in Fig.~\ref{fig:evidence_prior} were consolidated into a unified evidence-and-prior layer that became the main contextual input for the next stage of discussion. At this point, the workflow had not yet entered experimental planning, but it had already converted a broad and weakly structured optimization problem into a more explicit evidence landscape, including recurrent variables, major heterogeneities and evidence gaps, and sequence regions carrying potentially relevant prior signals. This integrated layer did not prescribe a final experimental plan in advance. Instead, it ensured that the subsequent project specification and planning stages began from a more structured and evidence-constrained understanding of the problem rather than from undirected exploration.

\subsection{Prior evidence was translated into executable wet-lab plans}

Building on the evidence-and-prior layer, the agent team reorganized the multivariable optimization problem around four experimentally resolvable questions: donor ranking under matched conditions, the choice between 55~$^\circ$C and 60~$^\circ$C, the practical value of extending reaction time to 4~h, and the contribution of donor-associated glycation beyond the common heating-and-ultrasound backbone. The main donor comparison therefore used lactose, glucose, galactose and mannose under a shared dry-state condition, while shorter time anchors and a restrained arabinose--ribose branch tested whether additional time or a faster-reacting donor offered greater value. Ultrasound, protein:sugar ratio, reaction mode and water activity were fixed to preserve comparability; broad wet-mode, ratio, pH and alternative-pretreatment screens were deferred.

The resulting planning draft contained 24 conditions. Before analysing outcomes, each planned condition was matched to the returned wet-lab records using donor, temperature, reaction time, ultrasound status, reaction mode, water activity and protein:sugar ratio. After incorporation of the supplemental ELISA record, all 24 planned conditions had directly attributable results and the original R01--R24 identifiers were retained for analysis (Table~\ref{tab:question_matrix}). The resulting matrix was organized into donor and temperature comparisons, no-ultrasound sentinels, shorter-time anchors, a limited pentose branch and heated no-sugar controls. The complete availability audit and source-record mapping are provided in Supplementary Tables~3 and~4.

\begin{table}[p]
\centering
\footnotesize
\setlength{\tabcolsep}{2.5pt}
\caption{
\textbf{Final question-structured wet-lab matrix with matched results.}
The initial 24-condition planning draft was retained as a complete 24-run matrix with matched results after supplemental ELISA records filled the previously unavailable conditions. The matrix was organized as five decision blocks rather than as a full factorial screen. All sugar-containing runs used $\beta$-LG variant A under dry-state glycation at $a_w$ 0.79 and a protein:sugar ratio of 1:2 (w/w); ultrasound status was varied only in the sentinel block and otherwise fixed as the default pretreatment. The table reports the comparison logic of the planned runs, not post hoc outcome-based selection.
}
\label{tab:question_matrix}
\begin{tabular}{@{}p{0.14\textwidth}p{0.10\textwidth}p{0.23\textwidth}p{0.24\textwidth}p{0.23\textwidth}@{}}
\toprule
\textbf{Decision block} & \textbf{Runs} & \textbf{Condition set} & \textbf{Question encoded} & \textbf{Comparator logic} \\
\midrule
Matched hexose bridge & R01--R08 (8 runs) & Lactose, glucose, galactose and mannose at 55~$^\circ$C or 60~$^\circ$C for 4~h, with ultrasound. & Which benchmark or challenger hexose performs best under a matched dry-state backbone, and does the preferred temperature lie at 55~$^\circ$C or 60~$^\circ$C? & Compare donors within each temperature arm; compare paired 55~$^\circ$C versus 60~$^\circ$C runs within each donor; interpret donor-containing runs against no-sugar process controls. \\
\addlinespace
No-ultrasound sentinels & R09--R12 (4 runs) & Lactose, glucose, galactose and mannose at 55~$^\circ$C for 4~h, without ultrasound. & Does the ultrasound-supported backbone provide a large independent gain under the matched 55~$^\circ$C hexose bridge? & Compare each no-ultrasound sentinel with its paired 55~$^\circ$C, 4~h, +US donor run while avoiding extrapolation to 60~$^\circ$C, pentose or long-time conditions. \\
\addlinespace
Benchmark and practical time anchors & R13--R16 (4 runs) & Lactose and glucose at 60~$^\circ$C for 2~h, and lactose and glucose at 55~$^\circ$C for 3~h, with ultrasound. & Does the 4~h bridge earn its additional time burden, or can shorter benchmark or practical anchors retain competitive value? & Compare 2~h or 3~h anchors with their corresponding 4~h donor branches; apply the pre-specified marginal-gain rule when added time produces less than a 10-percentage-point improvement. \\
\addlinespace
Restrained pentose branch & R17--R20 (4 runs) & Arabinose and ribose at 55~$^\circ$C for 2~h or 3~h, with ultrasound. & Should a fast-reacting pentose donor enter the main optimization path, and which observed short-window pentose condition is most credible? & Compare each pentose across 2~h and 3~h; compare arabinose and ribose at matched time points; judge pentose competitiveness against practical hexose anchors and process controls. \\
\addlinespace
Heated no-sugar process controls & R21--R24 (4 runs) & No donor at 55~$^\circ$C for 2, 3 or 4~h, and at 60~$^\circ$C for 4~h, with ultrasound. & How much signal reduction can be attributed to the common heating-and-ultrasound backbone rather than donor-associated glycation? & Serve as process-only references for short-time anchors and the matched 4~h sugar-containing temperature arms, preventing heating-only reduction from being overinterpreted as glycation-specific benefit. \\
\bottomrule
\end{tabular}
\begin{flushleft}
\footnotesize US, ultrasound pretreatment. The detailed run-sheet, availability audit, source records, sample labels, replicate identifiers and assay assignments are provided in Supplementary Tables~3--5. The corresponding machine-readable run-sheet and wet-lab data files are available in the project data repository.
\end{flushleft}
\end{table}

The matrix was designed for matched local comparisons rather than exhaustive optimization. All sugar-containing runs fixed $\beta$-LG variant A, a protein:sugar ratio of 1:2 (w/w), dry-state glycation and $a_w$ 0.79. The no-ultrasound sentinel block enabled a local 55~$^\circ$C, 4~h assessment of the ultrasound contribution for the four hexose donors, but it did not support extrapolation to 60~$^\circ$C, pentose donors or longer time windows. Reaction-time decisions followed a pre-specified marginal-benefit rule in which additional time was not preferred when it produced less than a 10-percentage-point improvement without a corresponding practical advantage.


\subsection{Iterative wet-lab results progressively identified high-value conditions}

The first wet-lab round was not intended to identify a global optimum directly across the full condition space, but to interrogate a limited set of high-information comparisons within the complete 24-run analysis matrix. Throughout this round, reduction values were calculated relative to the native untreated $\beta$-LG baseline. The returned data were interpreted through five comparison layers: donor-containing branches relative to the native untreated $\beta$-LG baseline, sugar-containing branches relative to heated no-sugar controls, 55~$^\circ$C versus 60~$^\circ$C within the matched 4~h bridge, paired no-ultrasound sentinels at 55~$^\circ$C and 4~h, and donor or short-window time-anchor comparisons under the dry-state backbone. These comparisons established several local but meaningful findings. First, donor-containing conditions substantially reduced IgG binding relative to the native untreated $\beta$-LG baseline. The no-sugar controls served as supportive process references, indicating that heating and ultrasound alone lowered the signal to a limited extent: no-sugar controls reached 15.47\%, 20.15\% and 25.94\% at 55~$^\circ$C for 2, 3 and 4~h, respectively, and 35.97\% at 60~$^\circ$C for 4~h, whereas donor-containing branches under the tested framework generally exceeded this range. Second, the gain from 60~$^\circ$C over 55~$^\circ$C was present but modest in the matched 4~h +US bridge. Mannose reached 68.15\% at 60~$^\circ$C compared with 65.40\% at 55~$^\circ$C, and the other hexoses showed similarly small differences, suggesting that the higher temperature imposed additional burden without materially changing the local branch ranking. Third, the no-ultrasound sentinels showed that ultrasound contributed only a modest local increment under the 55~$^\circ$C, 4~h hexose bridge: paired +US gains were 4.72, 1.81, 3.81 and 3.15 percentage points for lactose, glucose, galactose and mannose, respectively. Finally, pentose qualification signals were already visible in the first round, with ribose reaching 71.19\% at 2~h and 79.04\% at 3~h, and arabinose reaching 67.50\% and 74.98\% at the corresponding time points. Taken together, the first-round matrix successfully selected branches and narrowed the practical search region, but it did not yet establish a global optimum or support a general time-response law.

Because the first round completed branch selection without finalizing a processing route, the workflow then moved to a more focused second-round wet-lab block aimed specifically at marginal-benefit assessment. At this point, the key unresolved question was no longer which donor branch should remain under consideration, but whether additional reaction time within a promising branch still yielded sufficient practical gain to justify further continuation. This second-round block was therefore organized as a same-donor glucose time course under the fixed dry-state framework used at this stage of the workflow, with reductions again interpreted relative to the native untreated $\beta$-LG baseline. Unlike the first round, however, the purpose of this block was not donor ranking but within-branch slope estimation across time. Under this fixed framework, glucose reached 70.99\% reduction at 3~h, 78.75\% at 6~h, 82.50\% at 9~h, 84.33\% at 12~h, 87.91\% at 18~h, and 90.81\% at 24~h. This trajectory confirmed that time effects were real but clearly non-linear: higher reduction could be reached through longer extension, but not in a way that justified treating time as a uniformly strengthening axis or assuming that an additional hour would necessarily push the system across a target threshold. In this sense, the second round did not reinforce a simple longer-is-better logic; rather, it showed that pushing the hexose route deeper into the target region would require substantial additional time and process burden for progressively smaller gains.

On the basis of these two rounds of evidence-guided refinement, the workflow did not continue to extend the hexose branch indefinitely. Instead, it integrated the same-donor time-slope evidence with the earlier fast-donor signals and progressively shifted prioritization toward pentose donors as the more plausible route to strong IgG-binding reduction within a shorter practical window. In this narrowing process, the pentose route was no longer treated as a residual peripheral branch, but as a viable advancement path with continued potential for reduced immunoreactivity. Within that path, ribose progressively emerged as the more plausible pentose candidate. Final prioritization was not driven by reduction magnitude alone, but also considered reaction duration, process burden, evidence strength and practical experimental constraints. Accordingly, the workflow converged on a final evidence-weighted recommendation centred on ribose under a dry-state backbone with ultrasound pretreatment, corresponding to a 60~$^\circ$C, 4~h reaction at $a_w$ 0.79 and a protein:sugar ratio of 1:2 (w/w). The focused wet-lab validation of this final condition and its consistency with the pre-validation estimate are presented in the next section.


\subsection{Experimental validation of the final recommended condition}

Following the final evidence-weighted recommendation identified through iterative refinement, we next carried out focused wet-lab validation of the ribose-centred condition under the corresponding dry-state backbone. The final validated condition used ribose as the donor, with $\beta$-LG variant A at 10 mg mL$^{-1}$, under dry-state glycation at $a_w$ 0.79 and a protein:sugar ratio of 1:2 (w/w), following ultrasound pretreatment (400 W, 12 min) in 20 mM phosphate buffer at pH 7.3, with reaction at 60~$^\circ$C for 4~h. Relative to the native untreated $\beta$-LG baseline, this condition achieved an 85.6\% $\pm$ 1.7\% reduction in IgG binding (mean $\pm$ s.d.). This experimentally validated outcome was not only substantially higher than the strongest directly tested 4~h branch retained within the earlier workflow, namely mannose at 60~$^\circ$C for 4~h (68.15\%), but also confirmed that the final narrowing process had moved beyond the locally best directly tested hexose region toward a higher-value condition that remained experimentally supportable under focused validation.

The pre-validation recommendation included an evidence-weighted central estimate of 86\% IgG-binding reduction and a plausible range of 82--93\%. The experimentally observed value of 85.6\% $\pm$ 1.7\% (mean $\pm$ s.d.) was consistent with this range. The feedback and pre-validation recommendation record is summarized in Supplementary Table~6. Because the estimate and range were derived through evidence-weighted adjudication rather than a calibrated statistical prediction model, this agreement should be interpreted as support for the prioritization decision rather than as a formal measure of predictive accuracy. The practical value of the final ribose route was also not defined by reduction magnitude alone. Compared with the hexose route, which required substantial time extension to move further into the high-reduction region, the final ribose condition achieved strong IgG-binding reduction within a relatively compact 4~h window. Taken together, these results support the final validated ribose route as a balanced experimentally supported condition across IgG-binding reduction, reaction duration and process burden.

To further assess the validity and local robustness of the final ribose-centred condition, we conducted an additional confirmatory panel that was not part of the iterative planning loop. Instead, this auxiliary batch was designed as an independent check of whether the final 4~h ribose condition represented an isolated favourable point or belonged to a broader locally coherent high-performance region when calculated against the same native untreated $\beta$-LG reference. This panel covered an expanded donor-by-time map, including nine donors across 3, 6, 9, 12, 18, and 24~h under both ultrasound and no-ultrasound arms, together with an additional ribose 4~h validation point. Within this panel, ribose remained one of the strongest donors throughout the early-to-mid practical window: under ultrasound, ribose reached 79.04\% reduction at 3~h, 83.71\% at 6~h, 88.80\% at 9~h, and 91.19\% at 12~h, consistently ranking among the top-performing branches. These results indicate that the final ribose 4~h condition was not an isolated high-reduction observation, but part of a locally coherent high-value region. Moreover, the ribose branch retained a clear comparative advantage over several competing donors within the same practical window, while also showing a consistent directional benefit from ultrasound. Taken together, this auxiliary confirmatory panel further supports the reasonableness of the final ribose-centred recommendation by showing that it remained locally robust in neighbouring time points and extended donor comparisons, rather than depending on a single favourable observation.
The complete donor-by-time dataset, including ultrasound and no-ultrasound arms, is shown in Supplementary Fig.~1; the plotted values are provided as Source Data, and the raw and processed wet-lab workbooks are available in the project data repository.

\begin{figure}[!htbp]
  \centering
  \sffamily

  \begin{adjustbox}{max width=\linewidth, max height=0.82\textheight, keepaspectratio}
  \begin{minipage}{\linewidth}
    \centering

    \begin{minipage}[t]{0.48\linewidth}
      \centering
      \begin{overpic}[width=\linewidth]{Figure/Fig3/Figure_3a_v2.pdf}
        \put(1,53){{\sffamily\bfseries a}}
      \end{overpic}
    \end{minipage}
    \hfill
    \begin{minipage}[t]{0.48\linewidth}
      \centering
      \begin{overpic}[width=\linewidth]{Figure/Fig3/Figure_3b_v2.pdf}
        \put(1,53){{\sffamily\bfseries b}}
      \end{overpic}
    \end{minipage}

    \vspace{0.06in}

    \begin{overpic}[width=0.96\linewidth]{Figure/Fig3/Figure_3c_v2.pdf}
      \put(1,30){{\sffamily\bfseries c}}
    \end{overpic}

  \end{minipage}
  \end{adjustbox}

  \caption{
  \textbf{Wet-lab refinement and validation of the recommended glycation route.}
  \textbf{a}, Matched 4~h hexose bridge comparing lactose, glucose, galactose and mannose under 55~$^\circ$C +US, 60~$^\circ$C +US and 55~$^\circ$C no-US conditions.
  \textbf{b}, Short-window run-sheet anchors and no-sugar process controls measured under ultrasound-assisted dry-state conditions, comparing lactose, glucose, arabinose and ribose across the experimentally measured 55--60~$^\circ$C and 2--4~h conditions. All displayed points correspond to retained runs with condition-matched experimental records.
  \textbf{c}, Marginal-effect evidence and final validation. Left, the same-donor glucose time course at 55~$^\circ$C with ultrasound, which was provided to the agent team as a calibration source for non-linear time effects and diminishing marginal gains. For each pair of consecutive time points, the marginal gain was calculated as $[R(t_2)-R(t_1)]/(t_2-t_1)$, where $R(t)$ is the IgG-binding reduction at reaction time $t$. This quantity is expressed in percentage points per hour (pp h$^{-1}$). Right, the final recommended ribose condition at 60~$^\circ$C for 4~h, comparing the pre-validation estimate of 86\% (evidence-weighted plausible range, 82--93\%) with the observed IgG-binding reduction of 85.6\% $\pm$ 1.7\%. All reductions are expressed relative to the native untreated $\beta$-LG baseline.
  }
  \label{fig:wetlab_results}
\end{figure}


\subsection{Workflow accounting and division of labour in human--agent collaboration}

We used a seven-layer funnel to document how the experimental decision space changed across the workflow (Fig.~\ref{fig:workflow_efficiency}a). The starting envelope was derived from the structured literature JSONL evidence layer by harmonizing raw, paper-specific condition labels into semantically grouped literature-visible levels across six dimensions: donor branch, reaction-time window, temperature band, pretreatment family, protein:sugar-ratio band and environment-control state. This grouping merged synonymous donor labels and binned nearby or differently reported time, temperature, ratio and pH/RH/$a_w$ settings before estimating a conceptual upper bound of 1,123,584 combinations; the exact grouped-level counts and calculation are provided in Supplementary Methods~8 and Supplementary Table~8. This value describes the scale of the literature-derived variable space and was not treated as the number of experiments that would otherwise have been performed. Study scope and feasibility constraints first reduced this envelope to conditions relevant to the present $\beta$-LG system. Project specification then prioritized donor ranking, the 55--60~$^\circ$C comparison, reaction-time value, ultrasound sentinels and pentose advancement, which were translated into 24 planned conditions. The availability audit confirmed directly attributable results for all 24 planned runs, after which comparator-aware interpretation and focused refinement produced one pre-specified validation candidate. Thus, the funnel records a sequence of analytical and operational decisions rather than a proportional measure of experimental efficiency.

Interaction accounting showed a complementary division of labour between the human researcher and the agent roles (Fig.~\ref{fig:workflow_efficiency}b,c). Across the five analysed discussion phases, the human researcher contributed 2,146 words, primarily through objectives, feasibility constraints, tool-applicability checks and acceptability judgments, while the agent team produced 529,306 words of intermediate synthesis, criticism and planning. The Principal Investigator contributed 246,706 words, followed by the experimental design and evidence synthesis scientist (93,152), allergenicity and immunoassay interpretation scientist (85,806), Maillard reaction chemistry and process scientist (65,912), and scientific critic (37,730). Phase-normalized shares further showed that the Principal Investigator carried the largest textual contribution during project specification, tool selection and execution planning, whereas the experimental design role contributed more strongly during first-round interpretation. These counts describe the distribution of recorded interaction only; they are not direct measures of scientific quality, time saved or causal contribution.

The practical value of this division of labour was visible in the decisions retained by the workflow. The human researcher constrained the optimization target to conditions that were experimentally feasible and acceptable, while the evidence curator and domain roles organized prior observations and assay-specific considerations. The Principal Investigator integrated these inputs into testable comparisons, and the critic limited claims to the local evidence supported by matched comparators. Consequently, the glucose time course was interpreted as a marginal-benefit trade-off rather than as a reason for indefinite extension, donor prioritization considered process burden as well as IgG-binding reduction, and broad wet-mode, ratio, water-activity and pretreatment expansions were deferred when their expected information value did not justify the added confounding. The resulting accounting supports a traceable human--agent decision process, rather than a claim that textual volume alone demonstrates workflow superiority.
Complete phase-level and role-level interaction accounting is provided in Supplementary Table~7 and in the Source Data for Fig.~\ref{fig:workflow_efficiency}.

\begin{figure}[!htbp]
  \centering
  \sffamily

  \begin{adjustbox}{max width=\linewidth, max height=0.86\textheight, keepaspectratio}
  \begin{minipage}{\linewidth}
    \centering

    \includegraphics[width=0.98\linewidth]{Figure/Fig4/Figure_4a_search_funnel.pdf}

    \vspace{0.04in}

    \begin{minipage}[t]{0.49\linewidth}
      \centering
      \includegraphics[width=\linewidth]{Figure/Fig4/Figure_4b_phase_composition.pdf}
    \end{minipage}
    \hfill
    \begin{minipage}[t]{0.49\linewidth}
      \centering
      \includegraphics[width=\linewidth]{Figure/Fig4/Figure_4c_contributor_totals.pdf}
    \end{minipage}

  \end{minipage}
  \end{adjustbox}

  \caption{
  \textbf{Workflow accounting and division of labour in human--agent collaboration.}
  \textbf{a}, Seven-layer representation of the progressive narrowing process, from the literature-derived condition envelope to one pre-specified validation candidate. The initial envelope is a conceptual upper bound calculated from grouped literature-visible condition levels rather than a count of planned or feasible experiments; the execution-ready matrix reports 24 planned conditions and 24 runs with matched results after the supplemental availability audit. Intermediate widths are conceptual and are not proportional to candidate counts.
  \textbf{b}, Phase-normalized word shares across the five major discussion phases analysed in this study. Percentages are calculated within each phase and show the changing contribution of the human researcher and five agent roles. Human-researcher shares too small to fit inside their proportional segments are labelled immediately above the corresponding bars.
  \textbf{c}, Cumulative contribution across the same phases. The human researcher contributed 2,146 words and the agent team contributed 529,306 words in total. Counts include human-supplied prompts and questions and agent messages retained in the discussion and merged records; words were counted by whitespace-delimited splitting.
  }
  \label{fig:workflow_efficiency}
\end{figure}





